\begingroup
\setlength{\parskip}{3pt}
\setlength{\textfloatsep}{7pt plus 1pt minus 1pt}
\setlength{\abovecaptionskip}{3pt}
% Compact only the vertical gap before Section 3's run-in headings.
\makeatletter
\def\section{\@startsection{section}{1}{\z@}{-1.5ex plus -.3ex minus -.2ex}{1ex plus .2ex minus .2ex}{\large\sc\raggedright}}
\def\paragraph{\@startsection{paragraph}{4}{\z@}{5pt plus 1pt minus .5pt}{-1em}{\normalsize\bf}}
\makeatother
\section{Does visual generation help visual understanding?}
\label{sec:controlled_settings}
We first ask whether visual generation can improve visual understanding.
To isolate the effect of confounding factors, we construct paired I2I and I2T tasks that require solving the same task from the same input, differing only in whether the answer is produced as an image or as text.
This allows us to directly test whether supervision from an I2I objective transfers to the corresponding I2T task. This setup eliminates differences in annotation quality, image source, and other factors.
We study how the training recipe affects transfer, how I2I supervision reduces the need for I2T data, and whether the two objectives correspond at the instance level.
% Task construction, training details, and full recipe definitions are provided in Appendix~\ref{app:controlled}.
% \todo{if we add T2I should we test on some other models eg BLIP-3o as well?}
% Under this setup, we investigate three questions.
% First, can I2I supervision improve performance on the corresponding understanding task?
% Second, how should generation and understanding data be combined during training?
% Third, how does transfer depend on the relative data budgets of the two objectives?
\begin{figure}[t]
    \centering
    \includegraphics[width=\linewidth]{iclr2026/figures/combined_3subfigures_refined_v18.pdf}
    \caption{\textbf{Scaling I2I and I2T supervision.}
    (a) Controlled Jigsaw and Zoom-In tasks with image and text outputs.
    (b) I2T accuracy versus I2I training examples across recipes, with 1k I2T examples.
    (c) I2T accuracy versus I2T training examples for I2T-only and I2I $\rightarrow$ I2T training (100k I2I examples).
    Curves show means over three seeds. Shading indicates $\pm1$ SEM.}
    \label{fig:scaling}
\end{figure}
\paragraph{Controlled paired tasks.}
We adapt two tasks from VisGym~\citep{wang2026visgym}: Jigsaw and Zoom-In. 
In Jigsaw, patches of an image are shuffled, and the model must recover their spatial arrangement.
In Zoom-In, views of the same image at different zoom levels are shuffled, and the model must recover their correct order. 
For each input, the I2I objective produces the correctly ordered image, while the I2T objective predicts the same ordering as a permutation in text. 
The objectives require the same visual operation on the same input and differ only in output modality.
We instantiate this on BAGEL, a unified multimodal model based on a Mixture-of-Transformers (MoT) architecture, and then evaluate transfer exclusively on the paired I2T task.
% \paragraph{Text-to-Image}
% \todo{EXP: design and run experiments here, here I listed an idea; basically I2T be like (triangle, (3,1)); image is the rendered triangle at that coordinate; specifically synthetic scene rendering task:
% randomly generate different color, shape, coordination objects
% I2T: rendered image  →  canonical scene description, e.g.
% red square at (1,3), blue circle at (4,2).
% T2I: same canonical scene description → rendered image (refer to clever dataset?)}
% % Grounding: I2T is given image with bbox output coordiantas; T2I is giving a bbox coordinate draw a box on a canvas (?)
% We construct a synthetic scene task with objects of different shapes, colors, and locations. \todo{figure of illustration}
% For I2T, the model takes a rendered scene and outputs its structured description, including the attributes and coordinates of each object. For T2I, the model takes the same description and generates the corresponding scene.
% The two objectives therefore learn opposite directions of the same deterministic text-image mapping.
% experiment results, we found i2i help t2i does not; (why? any insight? i2i is learninng a delta, but for very toy t2i settings why does it not help)
\paragraph{Which training recipe transfers best?}
% show the scaling curve of I2I toy tasks under different recipes;
% show the curve udner different I2T data budgeste, enough I2I can appriximate I2T on toy settings
Under this setup, we study how I2I and I2T data should be combined during training. 
We compare six training recipes: \textbf{I2T-only}, \textbf{I2I $\rightarrow$ I2T}, \textbf{Mixed $\rightarrow$ I2T}, \textbf{Frozen I2I $\rightarrow$ Mixed}, \textbf{Mixed}, and \textbf{I2I $\rightarrow$ Mixed}.
We use $\rightarrow$ to denote sequential training stages.
\textbf{Mixed} denotes joint training with both I2I and I2T data in each batch.
\textbf{I2I} denotes training on visual generation data while updating parameters shared with the I2T objective.
\textbf{Frozen I2I} denotes visual generation training in which these shared parameters are frozen, and only visual generation-specific parameters are updated.
Architecture and training details are provided in Appendix~\ref{app:controlled-method}.
We fix the I2T budget at 1k examples with identical I2T training epochs across recipes and vary the amount of I2I data.
As shown in Figure~\ref{fig:scaling}(b), I2I $\rightarrow$ I2T and I2I $\rightarrow$ Mixed scale consistently as the I2I data increases.
Both recipes begin with an I2I-only training stage that updates parameters shared with I2T.
In contrast, freezing these weights during the I2I training stage or directly mixing the two objectives from the beginning leads to weaker or less stable gains.
We therefore use I2I $\rightarrow$ I2T as the default recipe in subsequent experiments.
\begin{findingbox}
An initial I2I training stage that updates parameters shared with the I2T objective provides a useful initialization for subsequent I2T learning.
\end{findingbox}
\paragraph{How much does I2I supervision reduce the need for I2T data?}
We next quantify how much I2I training can reduce the amount of I2T data needed to reach a given level of understanding performance.
We fix the I2I budget at 100k examples and compare I2I $\rightarrow$ I2T with I2T-only training at I2T budgets of 1k, 3k, 10k, and 30k examples.
As shown in Figure~\ref{fig:scaling}(c), I2I training improves performance across all four budgets, with larger gains at smaller I2T budgets.
On Zoom-In, 100k I2I examples followed by only 1k I2T examples achieve performance comparable to training on 10k I2T examples alone. 
On Jigsaw, 100k I2I examples followed by only 3k I2T examples achieve performance comparable to training on 10k I2T examples alone.
% 应该说的是可以substitute + suggest 当i2t data scarce但是i2i很多的时候可能很有用; 
Thus, evidence shows that I2I supervision can partially substitute for direct I2T supervision and may be particularly useful when understanding data is limited.
Training details are in Appendix~\ref{app:controlled-recipes}.
% \bs{Did we talk about scalability here? I think we mentioned that in the abstract.}
% \begin{findingbox}
% I2I supervision complements but does not replace direct I2T supervision,
% with the largest gains in low-I2T settings.
% \end{findingbox}
% Keep the table wrapped while allowing the paragraph to continue across pages.
\par\addvspace{3pt}
\begingroup
\setlength{\intextsep}{2pt}
\setlength{\columnsep}{5pt}
\begin{wraptable}[7]{r}{0.38\linewidth}
    \centering
    \small
    \setlength{\tabcolsep}{3pt}
    \setlength{\abovecaptionskip}{0pt}
    \makeatletter
    \def\@captype{table}
    \long\def\@makecaption#1#2{\vskip\abovecaptionskip{\small\raggedright #1: #2\par}\vskip\belowcaptionskip}
    \makeatother
    \caption{\textbf{I2I--I2T instance alignment.}}   \label{tab:controlled_instance_alignment}
    \begin{tabular}{@{}lcc@{}}
        \toprule
        & Jigsaw & Zoom-In \\
        \midrule
        I2I Acc. & 70.9 & 91.8 \\
        I2T Acc. & 81.3 & 90.6 \\
        $P(\mathrm{I2T}\checkmark\mid\mathrm{I2I}\checkmark)$ & 88.2 & 90.9 \\
        $P(\mathrm{I2T}\checkmark\mid\mathrm{I2I}\times)$ & 64.3 & 86.7 \\
        \bottomrule
    \end{tabular}
\end{wraptable}
\rightskip=0pt plus 1em
\noindent
\mbox{\textbf{Instance-level correspondence.}}\quad
Beyond aggregated transfer, we ask whether success on paired I2I and I2T tasks also corresponds at the level of individual examples.
We evaluate both I2I and I2T outputs from the same final I2I $\rightarrow$ I2T checkpoint for each of three training seeds, trained on 30k I2I and 1k I2T examples, using our 1k evaluation paired inputs.
We then compare I2T accuracy conditioned on whether the I2I prediction for the same input is correct.
As shown in Table~\ref{tab:controlled_instance_alignment}, I2T performance is higher on examples for which the corresponding I2I prediction is correct and lower on those that were wrong.
This indicates that the correspondence between paired I2I and I2T tasks extends to individual examples. More evaluation details are in Appendix~\ref{app:controlled-eval}
\par
\endgroup
\par\endgroup
